
\begin{figure*}[htbp]
    \small
    \begin{tcolorbox}[left=3pt,right=3pt,top=3pt,bottom=3pt,title=\textbf{Error Augmentation Prompt}]
    \textbf{\#\#\# Instruction}

    Suppose you are a careless student. I will provide you with a problem and the ground-truth solution. Your goal is to propose a new solution that contains specific errors. Specifically, you should:

    \begin{itemize}
        \item \textbf{Understand the Problem and Solution:} Carefully read the problem statement and the ground-truth solution provided. Ensure you fully understand the question and the solution. 
        
        \item \textbf{Introduce Specific Errors:}  Intentionally introduce a \textcolor{red}{\textbf{\{Error\_Type\}}} into the reasoning steps. 
        
        \textcolor{red}
        {\textbf{\{Error\_Description\}}}. Try to make these errors \textbf{plausible} and \textbf{reasonable} that even an intelligent student will make, not blatant or superficial. Don't mention the {error\_type} error explicitly. Your output should be as natural as a student's realistic reasoning process, in which he/she makes mistakes.

        \item \textbf{Continue Reasoning:} Continue reasoning after the error-injected step until you reach a wrong answer that is \textbf{inconsistent} with the final answer of the ground-truth solution. Leave the previous error as it is. Do NOT try to correct the error in the following steps.

        \item \textbf{Format Answer:} Give your final number in the format of 'The answer is [your answer]'. Stop at the wrong answer, do not try to correct it.

        \item \textbf{Explain the Error:} Explain why the proposed solution is incorrect and why it is a {error\_type} error. Your explanation should be wrapped in <explanation> tags.
    \end{itemize}

[Question]
{question}

[Ground-truth Solution] 
{gt\_solution}

[Error-Injected Solution]


\end{tcolorbox}
\caption{Error augmentation prompt for introducing flawed reasoning trajectories using the teacher model.}
\label{fig:introduce_error_prompt}
\end{figure*}